% !PRINTED 271
\weberpage{271}
%<<CONT
two conjugate pairs, and therefore $BB_1$ also has the form (13) and stands in the same relation to $AA_1$ as $B$ to $A$ and $B_1$ to $A_1$.

\gesperrt{From this it follows that the totality of the substitutions $B$ forms a group isomorphic with $\Lambda_p$, and this group we shall denote by $\Gamma_p$.}

If we put $\beta$ in place of $\nu\beta$ and denote by $\alpha', \beta'$ the quantities conjugate to $\alpha, \beta$, then we may also take the substitutions of the group $\Gamma_p$ in the simpler form:
\weq{14}{B = \begin{pmatrix}\alpha, & \beta\\ -\beta', & \alpha'\end{pmatrix}, \quad \alpha\alpha' + \beta\beta' = 1.}
In order to pass from one of these substitutions $B$ to the corresponding real substitution $A$, one first derives from (14) the corresponding substitution $A$:
\weq{15}{A = \begin{pmatrix}\alpha, & \nu^{-1}\beta\\ -\nu\beta', & \alpha'\end{pmatrix},}
and from it forms, according to (3) and (7):
\weq{16}{A = RAR^{-1}.}
In spite of the isomorphism, the group $\Gamma_p$ cannot be transformed into $L_p$ or $\Lambda_p$, at least not by substitutions whose numbers belong to the field $\mathfrak{E}_{2,p}$. To carry out the transformation, one would have to go into a higher field, in which \gesperrt{all} numbers of $\mathfrak{E}_{2,p}$ are squares.

\begin{center}
\S.\,70.\\
\gesperrt{Divisors of the group $L_p$ whose degree is}\\
\gesperrt{divisible by $p$.}
\end{center}

It is a problem of the greatest interest to determine all divisors of the group $L_p$. On the solution of this problem depends in what way the group $L_p$ can be represented by permutations of figures, and hence in what algebraic equations these groups can occur (\S.\,28, 2.).

The cyclic groups contained in $L_p$ we have already considered in the preceding sections. We have seen that the degree of a cyclic group is either equal to $p$ or a divisor of $1/2\,(p-1)$ or of $1/2\,(p+1)$
%CONT>>
